\documentclass{report}
\usepackage{amsmath,amssymb}
\setlength{\parindent}{0mm}
\setlength{\parskip}{1em}
\begin{document}
\begin{center}
\rule{6in}{1pt} \
{\large Dirk Abbeloos \\
{\bf A FMG half-space analysis for boundary control problems.}}

University of Leuven \\ Celestijnenlaan 200A \\ 3001 Heverlee
\\
{\tt dirk.abbeloos@cs.kuleuven.be}\\
Stefan Vandewalle\end{center}

We present a full multigrid (FMG) half-space analysis for one-shot
multigrid methods for time-dependent boundary control problems. The
control problem is written as an optimization problem consisting of a
linear parabolic PDE constraint and a quadratic tracking objective. More
precisely, we want to find the optimal control which steers the state of
the system as close as
possible to a given, predetermined state. In addition we regularize the
problem with a Tikhonov regularization term given by the squared L2-norm
of the control, which can be interpreted as the cost of the control.

Numerical experiments reveal strong divergence of classical multigrid
methods for small weights of the regularization term. Here classical is
interpreted as using well known modifications of standard multigrid
components near the boundary, e.g. restriction, prolongation, boundary
condition discretization, etc. For the interior equations well
established multigrid components were choosen, which illustrate that the
divergence is induced by the boundary control and not the interior
equations.

In order to analyse our multigrid method classical tools as for example
Fourier analysis are insufficient. They fail to incorporate the effect of
the boundary condition on the multigrid method. More precisely, they only
analyse the interior equations on an infinite or periodic grid, i.e. they
analyse the control problem without the boundary control. A more
complicated FMG half-space analysis allows us to analyse the effect of
the boundary. In our talk we present this FMG half-space analysis for a
selection of boundary control problems, as well as the insights obtained
from it. In addition we propose a modified multigrid method to solve the
given boundary control problem.


\end{document}
